\section{Discussion}

This study shows that neuron-detection performance cannot be interpreted safely
without preserving both annotation uncertainty and spatial identity. The frozen
pipeline recovered most known-positive occurrences, but the unresolved portion
did not reduce to one technical cause. Previously unmatched candidates included
provisional single-reviewer likely-neuron calls with appearances outside a
simple circular high-signal
template. Conversely, most local response improvements around known misses
crossed into existing identities. A final candidate score alone would obscure
both findings.

\subsection{Candidate discovery is useful before it is a precision estimate}

The \VenueReviewLikely{} definite or probable calls among
\VenueReviewedCandidates{} reviewed sites provide hypothesis-generating
evidence. They expand the set of appearances that future
annotation and proposal methods should accommodate, including small, weak, and
irregular structures. Their value does not depend on mislabeling the selected
fraction as precision. Independent second review can strengthen the individual
site decisions, while bounded exhaustive annotation is required to estimate how
often the detector is correct across the field.

This separation also clarifies the purpose of ranking. A review-value score may
rationally prioritize ambiguity, crowding, and possible failure modes. Such a
score can be negatively associated with obvious neuron calls and still be useful
for allocating expert effort. A neuron-probability score has a different target
and requires independent labels suitable for calibration.

\subsection{Localization and identity are different estimands}

Response-maximizing recentering appears attractive because it can increase
signal-to-background measurements. In crowded fields, however, the higher
response may belong to a neighbor. The \VenueIdentityCollision{} collision cases
demonstrate that a coordinate gain without an identity guard can overstate
recovery. The \VenueIdentityClear{} identity-clear cases expose a different
limitation: none was represented in the
nearby frozen proposal lists even at larger budgets. Improving those cases
requires proposal formation rather than reranking alone.

After independent review and bounded truth are completed, the result motivates
testing competition-aware candidate clusters that incorporate
distance, footprint overlap, peak separation, trace dependence, and split-versus-
two-source evidence. Outputs should support retain, merge, split, and abstain
decisions while preserving the original measurement. Evaluation should count
recovered identities, duplicate proposals, assignment stability, and abstention
coverage at a fixed review budget.

\subsection{Representation quality is task specific}

Raw fluorescence was strongest for frame-level retrieval, while spatial and
contextual features retained localized structure and ranking value. Registered
exact-truth simulations assigned the most robust proposal role under mechanism
shift to spatial context, while equal-weight combination was stronger only in
the original compact-budget factorial benchmark. A provisional legacy
two-frame audit remains an Appendix operator control rather than a main finding.
Together, these results discourage a single representation hierarchy.
Amplitude, context, stability, and source-separation outputs should be evaluated
for their intended roles rather than assumed to provide interchangeable evidence.

The exact-truth challenge suite reinforces the same boundary: crescent
morphology, close unequal-amplitude sources, deployable stopping, and parameter
stability remained limitations. These tests support mechanism-specific roles,
not a universal feature ordering or biological transfer.

\subsection{Limitations and next evidence}

All biological analyses come from one recording, so no population or
cross-recording generalization is claimed. Event windows and feature choices have
historical contact with this recording. Quiet intervals and displaced tissue are
not exhaustive negatives. The unmatched batch has one reviewer, and the
\VenueReviewUncertain{} uncertain calls remain unresolved. The
collision-versus-clear comparison contains only \VenueIdentityCollision{} and
\VenueIdentityClear{} cases. The simulation protocols provide known truth but
cannot reproduce the full acquisition or biological distribution.

The highest-value next evidence is therefore not a larger same-recording feature
search. It is independent adjudication of the \VenueReviewedCandidates{}
reviewed sites, exhaustive
annotation of a bounded field, and transfer of frozen rules to an independent
recording. A competition-aware resolver is a candidate design direction after
those truth and robustness gates and must be specified and frozen before any
evaluation used to support it. Motion and registration residuals should also be exported to
separate acquisition structure from neuronal observability.

The broader methodological lesson is that incomplete truth does not prevent
rigorous evaluation, but it changes the questions that can be answered. Known
positives support sensitivity and positive--unlabeled ranking. Matched controls
support operator and localization tests. Expert review generates hypotheses.
Only exhaustive or independent truth authorizes precision and generalization.
